\chapter{Protocol mechanisms}
There are currently four types of sensors supported by the protocol: a device, a
temperature sensor, a GPS sensor and a camera. Their functionality is as
follows.

A device sensor sends aperiodical updates about its status, that can be on or
off. A temperature sensor sends periodical updates (every second) about the
temperature. A GPS sensor is also aperiodical. Finally a camera sensor might
send no data at some time periods with bursts of binary data when the sensor
detects some motion.

As can be seen there are two fundamentally different classes of sensors: the ones
with periodical and often updates (like a camera and a temperature sensor), the
others are aperiodical and comparatively rarely update the data. The
specification is currently limited to the four aforementioned sensor types but it
makes the main difference between the two sensors classes. Mainly the difference
lies in reliability mechanisms.

The protocol discussed here is supposed to serve many different kind of
sensors. Basically the sensors are serving the state information of the sensor or some device the
sensor is monitoring. Among different users of this protocol it is likely that
we encounter some requirements for reliability of the protocol. In these cases
the protocol should be able to deliver the updates to clients even while some
messages might be lost. Such a situtation occurs for example with simple status
updates: if a sensor is in "OFF" state and next publishes its status to be "ON" but the message is lost the
client sees no change in the state of the sensor. The reliability with this kind
of messages is required as we see this kind of devices like state
machines and we also point out that in real deployment the client takes some action
based on the received status messages. Furthermore unreliable transfer of status
messages and other messages with similar requirements leads to situation where
one user might think the device is in "OFF" state while some others see it in
"ON" state. 

In some cases it is likely that a loss of some updates can be
tolerated. Such is the case for example with streaming data i.e. updates that
are soon outdated. In these cases it is likely that the reliability guarantee
will only lead to wasted resources.

To manage these requirements the client can specify in a \texttt{subscribe}
message a specific flag requiring the updates to be delivered in reliable
manner. When this reliability parameter is used the server requires that client
acknowledges updates. If no acknowledgment is received within a given
timeout the server will retransmit the update message to the client. However the
protocol does not guarantee that all messages will be delivered to clients when this
optional reliability is used. If a new update is received by the server from a
sensor and there are still some active retransmissions to clients from this same
sensor the newer update will replace the old message. We refer this behaviour as
\textit{weak-reliability}. We argue that for IoT protocol like this the
weak-reliability is enough as many kind of applications need only the most
recent information. Applications requiring more strict reliability guarantees
should fall back using plain old TCP connections. 

% XXX: Reset to zero

In order to avoid reordering of messages all subscriptions have a logical time value which
is incremented with each message. A client receiving an update from a sensor
with smaller time value than the previous update from this same sensor must
discard the message as old. When the timer value of the last received message is
high enough the receiver should accept also timer values with very small timer
values. This is because the timer value is basically an unsigned integer which
will overflow and reset to zero when the value grows big enough.

% TODO add about keep-alive's here
The server keeps a list of subscribed clients per sensor and the list used when an
update must be delivered to subscribed clients. However because the publish-subscribe
protocol uses UDP on transport layer the server has no information whether or not
the updates are reaching the client when the connection is specified as unreliable.
In order to avoid server to still handle "dead clients" in its receiver lists a client
must eventually send \texttt{keep-alive} messages to server. The keep-alive timeout
is a protocol parameter used by both servers and clients. The server should use a
a client keep-alive timeout multiplied by some constant to avoid unnecessary client removals.
If a client is removed from the server it no longer receives updates from the server and it must
resubscribe with the normal subscribe procedure.

The recommended value for the client keep-alive timeout is 5 seconds. And the
recommended value for the server constant is 5. Thus if no keep-alive is
received within 25 seconds, the client is pruned from the client list.



